% !TeX root = ../../../../../main.tex
% chapters/sections/chapter4/subsections/parameters/monetary_optimization.tex

\subsection{自国金融政策の最適パラメータの決定とゲーム的状況の回避}
\label{sec:parameters_optimization_and_game_theory}
本稿では\ref{chap:chapter5}章において様々な自国金融政策のパフォーマンスを比較する。
その際、各金融政策が本来持つ性質がもたらす結果の差異を比較するために、
パラメータの設定方法は統一された基準に従う必要がある。
もし個々の金融政策について恣意的にパラメータを設定するなど統一を欠いた場合、
パラメータの設定方法の違いによってシミュレーション結果が左右されてしまう可能性があるためである。
そこで本稿では、すべての金融政策について妥当性をもった比較をおこなうための統一基準として
「\(\beta\)ショックのシミュレーションにおいて自国の厚生を最大化する」という条件下で
それぞれの最適パラメータを探索し決定する。
しかし、2国モデルにおいて自国がパラメータを動かしながら最適化を試みると、
通常はその影響が他国のマクロ経済変数にも波及し、
他国にとっても最適な政策パラメータが変化してしまうというゲーム理論的な相互依存状況が生じる。
本稿ではこの複雑な状況を回避するため、外国の金融政策を生産者物価指数を用いたインフレ目標に固定し、
その政策パラメータも標準的な値に据え置くという設定を採用する。
この非対称な設定は、\ref{chap:chapter5}章で詳述する本モデル特有の「遮断効果」によって正当化される。
本モデルの構造上、自国の金融政策の変更は外国の生産者物価指数などの主要な外国関連変数に影響を与えない。
そのため、自国がどのような政策パラメータを選択したとしても外国側の最適応答が変化することはなく、
外国の政策を固定したまま自国のパラメータ最適化のみに集中することが可能となる。
表\ref{tab:chapter4_parameters_list}の下段に示された日本（\(H\)）の各種金融政策のパラメータは、
以上の前提の下で探索された最適値であり、
これらを用いておこなったシミュレーション結果が\ref{chap:chapter5}章で示される。